\documentclass{article}

\usepackage[margin=1in]{geometry}
\usepackage{booktabs}
\usepackage{graphicx}
\usepackage{hyperref}

\title{Speculative Any-Order Adaptive Editing for Faster Diffusion Language Models}
\author{Pranav Sitaraman}
\date{}

\begin{document}
\maketitle

\begin{abstract}
We study whether a block-diffusion language model can be retrofitted with a
speculative semi-any-order policy that trades verifier compute for higher
throughput while preserving answer quality. AOAE uses a cheap hard-routed
drafter, a stronger verifier, and a scalar learned policy over unmask/remask
decisions. The final experiments train with GRPO on GSM8K and evaluate on
GSM8K, MATH-500, and HumanEval. Cache/access rewards are disabled in the final
run; the reported frontier isolates speculative editing and semi-any-order
policy expressiveness.
\end{abstract}

\section{Claims}

The paper tests three claims. First, speculative AOAE gives direct control over
the LLaDA2.1 speed/accuracy tradeoff. Second, semi-any-order editing increases
expressiveness relative to strict block decoding, although it may spend more
verifier compute. Third, GRPO moves the built frontier mostly upward in
accuracy and somewhat rightward in throughput by teaching the scalar policy to
choose better unmask/remask actions.

\section{Method}

The final implementation uses LLaDA2.1-mini with a hard-routed auxiliary drafter
and soft verifier. The policy is Phase-A/V2 scalar-only: it trains the
unmask/remask heads and excludes hidden-state residual features, stable-cache
execution, cache reward terms, and access reward terms. The canonical config is
\texttt{configs/paper.yaml}; eval configs are \texttt{configs/eval\_gsm8k.yaml},
\texttt{configs/eval\_math500.yaml}, and \texttt{configs/eval\_humaneval.yaml}.

\section{Experiments}

Training uses GSM8K-style math prompts through the warmstart then GRPO path.
Evaluation reports in-distribution GSM8K and out-of-distribution MATH-500 and
HumanEval. The default trained checkpoint is the V4 quality-balanced scalar
policy at \texttt{outputs/v4\_grpo\_quality\_balanced/policy\_best.pt} when
available. The no-train baseline uses the same eval configs with
\texttt{--checkpoint none}.

\section{Results}

Final artifacts are generated under \texttt{outputs/paper\_final/}. The paper
build uses the exact manifest in \texttt{outputs/paper\_final/run\_manifest.json}
and the aggregate table in
\texttt{outputs/paper\_final/aggregate\_comparison\_table.md}.

\begin{figure}[h]
\centering
\IfFileExists{../outputs/paper_final/gsm8k/trained/tps_accuracy_pareto.png}
{\includegraphics[width=0.72\linewidth]{../outputs/paper_final/gsm8k/trained/tps_accuracy_pareto.png}}
{\fbox{\parbox{0.68\linewidth}{GSM8K Pareto plot generated by
\texttt{bash reproduce.sh --workflow paper}}}}
\caption{GSM8K speed/accuracy frontier.}
\end{figure}

\begin{figure}[h]
\centering
\IfFileExists{../outputs/paper_final/math500/trained/tps_accuracy_pareto.png}
{\includegraphics[width=0.72\linewidth]{../outputs/paper_final/math500/trained/tps_accuracy_pareto.png}}
{\fbox{\parbox{0.68\linewidth}{MATH-500 Pareto plot generated by
\texttt{bash reproduce.sh --workflow paper}}}}
\caption{MATH-500 out-of-distribution frontier.}
\end{figure}

\begin{table}[h]
\centering
\begin{tabular}{llrr}
\toprule
Run & Dataset & Accuracy & TPS \\
\midrule
LLaDA2.1 speed & GSM8K & 74.60 & 72.6 \\
LLaDA2.1 quality & GSM8K & 77.26 & 64.7 \\
V4 QBal best, tau0.5 lossy & GSM8K & 79.61 & 72.6 \\
LLaDA2.1 quality & MATH-500 & 41.40 & 91.5 \\
V4 QBal best, tau0.5 lossy & MATH-500 & 36.40 & 112.5 \\
\bottomrule
\end{tabular}
\caption{Preliminary full-set results preserved from the final digest. The
camera-ready table should be regenerated from
\texttt{outputs/paper\_final/aggregate\_comparison\_table.csv}.}
\end{table}

\section{Limitations and Future Work}

AOAE is not yet a theoretically full any-order diffusion decoder. The evidence
supports semi-any-order retrofitting as a promising direction, analogous in
spirit to LLaDA2.1's distillation of autoregressive behavior into block
diffusion. Future work should train policies and model objectives that make
global any-order decoding native rather than retrofitted.

\bibliographystyle{plain}
\bibliography{main}

\end{document}
